Este reporte está organizado alrededor de una idea central: no basta con construir grafos Mapper y observar que tienen estructura; también es necesario evaluar qué tipo de evidencia sostiene esa estructura. En este proyecto, Mapper se usa como una herramienta exploratoria para revelar regiones candidatas en el espacio de propiedades de exoplanetas. Después, esas regiones se auditan desde tres dimensiones: su dependencia de imputación, su posible interpretación física u orbital y su asociación con variables de sesgo observacional.

La lectura general del proyecto es que la estructura más informativa aparece en el espacio orbital. Sin embargo, esa estructura no puede tomarse directamente como una señal astrofísica pura, porque la auditoría por método de descubrimiento muestra una asociación estadísticamente significativa entre la topología del grafo y la forma en que los planetas fueron detectados. Por eso, el resultado más importante no es una clasificación cerrada de exoplanetas, sino una clasificación crítica de las regiones observadas: algunas parecen físicamente interpretables, otras parecen dominadas por sesgo observacional, algunas son mixtas y otras deben considerarse débiles por su dependencia de imputación o falta de evidencia suficiente.

\begin{table}[H]
\centering
\small
\begin{tabular}{p{0.43\textwidth}p{0.43\textwidth}}
\toprule
\textbf{Elemento evaluado} & \textbf{Resultado principal} \\
\midrule
Mappers reconciliados & 21 grafos completos \\
Diseño de comparación & 7 espacios de variables $\times$ 3 lentes \\
Cubierta usada & \texttt{cubes10\_overlap0p35} \\
Capa principal de interpretación & \texttt{pca2} \\
Grafo candidato principal & \texttt{orbital\_pca2\_cubes10\_overlap0p35} \\
Nodos del grafo orbital & 124 \\
Aristas del grafo orbital & 196 \\
Ciclos del grafo orbital & $\beta_1=96$ \\
Imputación media nodal orbital & 0.011 \\
Imputación media nodal térmica & 0.588 \\
Pureza observada por método en orbital & 0.842 \\
NMI entre nodo y \texttt{discoverymethod} & 0.254 \\
$p$ empírico por permutación & 0.001 \\
Lectura principal & Señal orbital informativa con sesgo observacional significativo \\
\bottomrule
\end{tabular}
\caption{Resumen ejecutivo de los hallazgos medibles principales del proyecto.}
\label{tab:executive_summary}
\end{table}

La Tabla~\ref{tab:executive_summary} resume la lógica de evidencia del proyecto. El Mapper orbital es el resultado más prometedor porque combina alta complejidad topológica con baja imputación. En cambio, el Mapper térmico muestra que una estructura compleja no necesariamente es más confiable si depende fuertemente de valores imputados. Esta comparación es importante porque obliga a distinguir entre una estructura visualmente rica y una estructura científicamente defendible.

La auditoría de sesgo observacional refuerza esta cautela. El hecho de que el grafo orbital tenga una asociación significativa con \texttt{discoverymethod} indica que parte de la topología observada puede estar relacionada con el proceso de recolección del catálogo. Esto no invalida el resultado orbital, pero sí cambia su interpretación: el grafo no debe leerse como una clasificación directa de arquitecturas planetarias, sino como una estructura donde se mezclan señal orbital, sesgo de descubrimiento y efectos del preprocesamiento.

En consecuencia, el reporte adopta una postura conservadora. Las regiones de Mapper se interpretan como candidatas para inspección astrofísica, no como clases finales. Una región será más confiable si combina baja imputación, coherencia física u orbital y baja asociación con método de descubrimiento. Una región será más sospechosa si está dominada por un método, una facilidad de descubrimiento o una época específica. Esta distinción permite que el análisis sea más fuerte: no solo muestra que hay estructura, sino que evalúa qué tan interpretable es esa estructura.

